\chapter{Архитектура целиком}
\label{ch:architecture}

Прежде чем разбирать механизмы по отдельности, полезно увидеть систему целиком:
из чего она состоит, что по ней течёт и как компоненты связаны. Эта глава — карта.
Всё, что здесь названо, ниже определяется формально (глава~\ref{ch:definitions}) и
разбирается подробно; сейчас задача — общий взгляд.

\section{Что мы вообще строим}

Рин — локальный AI-компаньон. У неё должна быть \emph{память}, которая ведёт себя не
как журнал сообщений, а как память живого существа: связная, ассоциативная,
меняющаяся во времени, способная ошибаться и исправляться, окрашенная настроением и
характером. Центральная ставка проекта: такую память естественно представить
\emph{графом}, а всё «поведение памяти» — процессами (операторами) над этим графом.

\begin{principle}[Единый субстрат]
Память Рин — это \textbf{один типизированный темпоральный граф с метаданными на
узлах и рёбрах}. Retrieval, вывод (fusion), консолидация, забывание, ночная работа
(сон), настроение и слой личности — не отдельные подсистемы со своими базами, а
\emph{операторы}, читающие и пишущие одни и те же поля одного графа.
\end{principle}

\begin{pitfall}[«Один субстрат» — дисциплина, а не «всё есть граф»]
Тезис легко выродить в догму, и тогда в граф начнут запихивать всё подряд:
runtime-очереди, streaming-флаги, транзиентные тики настроения, состояние
планировщика. Это превратит граф в свалку «всего, что случилось». Граница проходит
чётко: \textbf{в графе живёт долговременное содержание и его эпистемические
обоснования} — воспоминания, факты, убеждения, связи, провенанс. Остальное — рядом,
и это \emph{не} граф:
\[
\text{граф} = \text{память};\quad
\text{event log} = \text{история};\quad
\text{телеметрия} = \text{самонаблюдение};
\]
\[
\text{scheduler state} = \text{управление};\quad
\text{prompt context} = \text{временная сборка для LLM}.
\]
Они связаны и обмениваются данными, но это разные хранилища с разными законами
жизни (главы~\ref{ch:runtime},~\ref{ch:mood}). Настроение — сцепленный сбоку модуль;
самотелеметрия метакогниции — кольцевой буфер; планировщик — конечный автомат.
Правило простое: если данное \emph{транзиентно} или это \emph{управление/история} —
ему не место в графе памяти.
\end{pitfall}

\section{Слои системы}

Система распадается на четыре слоя. Снизу вверх:

\begin{enumerate}
\item \textbf{Субстрат} (глава~\ref{ch:substrate}) — сам граф: узлы-воспоминания и
типизированные знаковые рёбра, с полями \texttt{confidence}, \texttt{valence},
\texttt{provenance}, временной слой, \texttt{activation}, флаги \texttt{need\_to\_*}.
Его начальное наполнение — \emph{эмбриональный граф}
(глава~\ref{ch:embryonic}) из готовых баз здравого смысла.
\item \textbf{Операторы} — процессы над субстратом:
\emph{обход} роем случайных блуждающих агентов (глава~\ref{ch:walker}),
\emph{вывод и консолидация} (глава~\ref{ch:fusion}),
\emph{пропагация уверенности} (глава~\ref{ch:propagation}),
\emph{забывание} (глава~\ref{ch:timescales}),
\emph{ночная работа/сон} (глава~\ref{ch:sleep}).
\item \textbf{Состояния} — сцепленные с графом, но отдельные:
\emph{настроение} (глава~\ref{ch:mood}) и производные эпистемические статусы
(глава~\ref{ch:epistemics}).
\item \textbf{Выражение} — то, что превращает внутреннее в реплику:
\emph{личность} (глава~\ref{ch:personality}) и \emph{голос/языковая модель}
(глава~\ref{ch:voice}).
\end{enumerate}

\begin{center}
\begin{tikzpicture}[font=\small,>=Stealth,
  layer/.style={draw,rounded corners,minimum width=82mm,align=center,inner sep=6pt,fill=gray!5},
  lbl/.style={font=\bfseries}]
  \node[layer] (expr) {{\bfseries Выражение}\\[1pt]
     \footnotesize личность \quad$\cdot$\quad Rinness / языковая модель};
  \node[layer,below=7mm of expr] (state) {{\bfseries Состояния}\\[1pt]
     \footnotesize настроение (mood) \quad$\cdot$\quad уверенность / статусы};
  \node[layer,below=7mm of state] (op) {{\bfseries Операторы}\\[1pt]
     \footnotesize обход (walker) $\cdot$ fusion $\cdot$ консолидация\\
     \footnotesize пропагация $\cdot$ сон / скраб};
  \node[layer,below=7mm of op] (sub) {{\bfseries Субстрат}\\[1pt]
     \footnotesize граф: узлы + знаковые рёбра + метаданные\\
     \footnotesize эмбриональный граф};
  \draw[<->,thick] (expr)--(state);
  \draw[<->,thick] (state)--(op);
  \draw[<->,thick] (op)--(sub);
\end{tikzpicture}
\captionof{figure}{Четыре слоя системы. Стрелки двусторонние: настроение смещает
обход, обход меняет граф, граф питает настроение, и так далее.}
\label{fig:layers}
\end{center}

\section{Что течёт по системе}

Полезно различать четыре «потока», все — над одним графом:

\begin{itemize}
\item \textbf{Активация} — транзиентная «подсветка» узлов при обходе; быстро гаснет.
\item \textbf{Уверенность} — медленно меняющееся поле $\conf$; её пропагация и
ревизия — эпистемическое ядро.
\item \textbf{Валентность и настроение} — аффективная окраска узлов и сцепленный
сбоку вектор настроения.
\item \textbf{Структурные изменения} — рождение узлов (fusion), укрепление/распад
рёбер (консолидация/забывание): граф эволюционирует.
\end{itemize}

\section{Жизнь одного воспоминания}

Сквозной проход, чтобы увидеть, как слои работают вместе. Пусть на стриме зритель
пишет: «Рин, я опять слил ранг в доте».

\begin{workedbox}[От реплики до реплики]
\textbf{1. Восприятие.} Событие кодируется как узел \texttt{ev}: «зритель $Z$
проиграл». Оценивается его значимость (глава~\ref{ch:personality}); заметный сдвиг
настроения делает событие «важным».

\textbf{2. Обход (retrieval).} Рой блуждающих агентов стартует от \texttt{ev},
смещённый текущим настроением Рин (озорным) и запросом. Всплывает окрестность:
«проигрыш $\to$ досада/смех», «$Z$ раньше хвастался скиллом», «дота».

\textbf{3. Вывод (fusion).} Совместная активация «$Z$ хвастался» и «$Z$ проиграл»
рождает \emph{гипотезу}-узел «повод беззлобно подколоть $Z$» — провизорно, с низкой
уверенностью и провенансом $\{$хвастался, проиграл$\}$.

\textbf{4. Выражение.} Личность (участие, колкая нежность) и настроение выбирают
\emph{тон}; Rinness-модель \emph{композирует} реплику из всплывших ингредиентов —
не берёт готовую фразу из графа. На выходе: «Оёй, ну куда без слитого ранга. Иди
чаёк попей, чемпион».

\textbf{5. Запись.} Реплика и реакция $Z$ пишутся провизорно. Если $Z$ смеётся —
уверенность ребра «подкол работает с $Z$» растёт; если обиделся — гаснет, и
обновляется модель отношения с $Z$.

\textbf{6. Ночью.} Во сне (глава~\ref{ch:sleep}) телепортирующий обходчик случайно
переинспектирует память, консолидирует частые эпизоды («$Z$ — частый гость, любит
доту»), проверяет противоречия, иногда рождает неожиданные ассоциации.
\end{workedbox}

Каждый шаг — отдельная глава. Дальше мы разбираем их по одному, но держите в голове
именно этот сквозной цикл: \emph{восприятие $\to$ обход $\to$ вывод $\to$ выражение
$\to$ запись $\to$ ночная переработка}.

\section{Карта чтения}

Читателю-новичку достаточно глав~\ref{ch:architecture} (эта),
\ref{ch:definitions} (словарь) и сквозного примера выше, чтобы понимать разговоры
команды. Дальше главы можно читать по интересу: субстрат и эмбриональный граф — для
тех, кто будет строить хранилище; обход, fusion, сон — для тех, кто пишет
операторы; настроение, личность, голос — для тех, кто отвечает за характер;
эксперименты (глава~\ref{ch:experiment}) — для всех, потому что проверять придётся
всё.
